\section{Introduction}

We trained ResNet-50 on Waterbirds and found that --- without any group labels --- the curvature of the loss surface around each training sample reveals minority group membership with AUROC 0.85--0.90 in 4/5 random seeds (mean AUROC 0.89 across the 4 passing seeds). But when we investigated \emph{why}, the expected mechanism was absent: by the time the curvature signal peaks, minority samples are classified just as confidently as majority samples. The method works --- but not for the reason we expected.

This puzzle sits at the intersection of two long-standing problems in robust machine learning. The first is spurious correlations: models trained with empirical risk minimization (ERM) on real-world datasets learn to exploit incidental statistical associations between input features and labels. On Waterbirds \cite{sagawa2019distributionally}, a ResNet-50 trained with ERM achieves 97\% average accuracy but only 72\% worst-group accuracy --- minority groups (e.g., landbirds on water backgrounds) suffer because the model relies on background rather than bird shape. The second problem is annotation burden: the most effective mitigation methods, such as deep feature reweighting (DFR) \cite{kirichenko2022last}, require a balanced held-out validation set with explicit group labels --- a resource unavailable in many practical settings.

Annotation-free alternatives --- Just Train Twice (JTT) \cite{liu2021just}, SELF \cite{labonte2023towards}, EVaLS --- proxy minority membership through first-order signals: which samples are misclassified, or which have the highest loss. These methods improve over ERM without group labels, but they operate on signals that are fundamentally insensitive to \emph{how} the model encodes spurious features in its loss landscape geometry.

We ask: does the second-order geometry of the loss surface carry distinct information about minority group membership? Specifically, we measure the per-sample Hessian trace of the last fully-connected layer --- computed via $K{=}50$ Hutchinson estimation using \texttt{torch.func} \texttt{vmap+vjp} --- at six training checkpoints ($t\in\{0,1,5,10,20,50\}$) over five random seeds, and ask whether this trajectory discriminates minority from majority samples.

The measurement protocol is motivated by two theoretical results. \cite{labonte2024spectral} show that minority group covariance matrices have larger spectral norm than majority groups --- a group-level prediction of systematic feature-norm inequality ($\|x_i\|^2$) that would elevate per-sample Hessian traces for minority samples. \cite{labonte2026phase} prove that SGD on spurious data learns the spurious feature first (Phase I) before the core feature (Phase II), predicting a transient differential in loss landscape curvature between majority (which gain confidence rapidly via the spurious shortcut) and minority (which do not).

Our existence experiment (H-E3) confirms the signal: 4/5 seeds achieve AUROC$\geq$0.85 for minority membership prediction at $t^*$ (the epoch maximizing the trace ratio $R(t) = \overline{\text{trace}}_{\text{min}} / \overline{\text{trace}}_{\text{maj}}$), with epoch-0 AUROC$<$0.70 in all seeds --- confirming that ERM training creates the signal, not the pretrained ImageNet initialization. The Hutchinson estimator is stable (coefficient of variation CV$<$3\% for all seeds), establishing $K{=}50$ as sufficient for the last-fc layer of ResNet-50.

Our mechanism experiment (H-M1) produces the puzzle. The original mechanistic explanation --- that minority samples remain near the decision boundary ($p_i\in[0.3,0.7]$) throughout Phase I, keeping the confidence-entropy term $p_i(1-p_i)$ in $\text{Tr}(H_i^{fc}) \propto \|x_i\|^2 p_i(1-p_i)$ large --- is directly falsified. At $t^*$, minority training-set confidence is 0.9678--0.9999 across all seeds: fully saturated, not boundary-straddling. The confidence channel is inactive. A transient differential does exist at epoch 1 (seed 1: $p_{\min}=0.827$ vs $p_{\text{maj}}=0.972$), consistent with Phase I theory, but it disappears before $t^*$.

The trace signal at $t^*$ must therefore arise from the $\|x_i\|^2$ feature-norm channel --- systematic differences in the penultimate-layer representation norms between minority and majority samples, consistent with \cite{labonte2024spectral}'s spectral imbalance finding. This is an unverified but testable candidate mechanism that we identify as the primary direction for future mechanistic investigation.

Our contributions are:

\begin{enumerate}
\item \textbf{Empirical existence result:} First experimental evidence that per-sample last-fc Hessian trace ($K{=}50$ Hutchinson via \texttt{vmap+vjp}) achieves AUROC$\geq$0.85 for minority membership detection in ERM-trained ResNet-50 on Waterbirds across 5 random seeds, with epoch-0 control confirming ERM emergence (Section~\ref{sec:results}).

\item \textbf{Principled falsification of the confidence-differential mechanism:} Direct refutation of the sustained decision-boundary hypothesis ($p_{\text{minority}}\in[0.3,0.7]$ at $t^*$) on training-set data, identifying the feature-norm channel ($\|x_i\|^2$) as the likely driver (Section~\ref{sec:discussion}).

\item \textbf{Practical confirmation:} $K{=}50$ Hutchinson is the minimum stable setting for last-fc ResNet-50 (CV$<$3\%), with the epoch-0 control serving as a reliable pretrained-artifact detector (Section~\ref{sec:experiments}).
\end{enumerate}

The remainder of the paper is organized as follows. Section~\ref{sec:related} reviews related work. Section~\ref{sec:method} describes our methodology. Section~\ref{sec:experiments} details the experimental setup. Section~\ref{sec:results} presents results. Section~\ref{sec:discussion} discusses findings, limitations, and the mechanism puzzle. Section~\ref{sec:conclusion} concludes.
